% STATUS: written 2026-11-05 (Wang 2016 read in full; Liu-Liu 2021 abstract read)
% STATUS: extended 2026-11-06 (Guan and Zhang 2024, Di Bartolo and La Rosa 2026, Di Bartolo, Di Fatta and La Rosa 2025 (both papers) read in full)
% STATUS: extended 2026-11-07 (Cheraghpour and Jafari 2024 read in full from the
% library PDF)
% Sources: Wang 2016 read in full (pp. 1952-1959). Liu-Liu 2021: only the abstract and the
% references extract from the local copy (see excerpts/Liu2021_jordanbid.txt). Bresar 1995 is
% atlas-pending P0005 (local copy image-only); the generalized-biderivation notion is quoted
% through Yazarli 2017 (p. 490, attribution) and through the Wang 2016 triangular-ring theory
% that builds on Bresar-Martindale-Miers 1993 via the notion of an inner biderivation.
% Guan and Zhang 2024 (pp. 1--4, 13); Di Bartolo and La Rosa 2026 (pp. 1--5); Di Bartolo, Di Fatta
% and La Rosa 2025 Leibniz (pp. 1, 6--8, 10); the same authors 2025 Super (pp. 1--2, 5, 10);
% Cheraghpour and Jafari 2024 (pp. 1--10 of the arXiv version).
\section{Tri-derivations, generalized biderivations and Jordan biderivations}\label{sec:higher-multi:tri-generalized}

Sections~\ref{sec:higher-multi:biderivations}
to~\ref{sec:higher-multi:symmetric-maps} studied a biderivation in prime and
semiprime rings. Brešar extends the notion to a generalized biderivation and
develops its properties \cite[reference~5]{Bresar1995GBid}. This section looks
at what a biderivation can be on a structured ring
\cite[p.~1952]{Wang2016Bideriv}.

Two special forms carry the theory \cite[p.~1952]{Wang2016Bideriv}.

\begin{definition}[{\cite[p.~1952]{Wang2016Bideriv}}]\label{def:higher:inner-extremal}
Let $R$ be a noncommutative ring. A biderivation $\phi\colon R\times R\to R$
is \emph{inner} when $\phi(x,y)=\lambda[x,y]$ for all $x,y\in R$, where
$\lambda\in Z(R)$. A biderivation $\phi\colon R\times R\to R$ is
\emph{extremal} when $\phi(x,y)=[x,[y,a]]$ for all $x,y\in R$, where $a\in R$
with $a\notin Z(R)$ and $[[R,R],a]=0$.
\end{definition}

On a triangular ring every biderivation is built from these two
\cite[p.~1955]{Wang2016Bideriv}. A unital ring $R$ with a nontrivial
idempotent $e$ is \emph{triangular} when $f=1-e$, $eRf$ is a faithful
$(eRe,fRf)$-bimodule and $fRe=0$ (Theorem~\ref{thm:higher:triangularbid}) \cite[p.~1953]{Wang2016Bideriv}.

\begin{theorem}[{\cite[Theorem~2.1, p.~1955]{Wang2016Bideriv}}]\label{thm:higher:triangularbid}
Let $R$ be a triangular ring. If the following conditions hold:
\begin{enumerate}
\item $C(fQ_{ml}(R)f,fRf)=C(R)f$,
\item either $eRC(R)e$ or $fRC(R)f$ does not contain nonzero central ideals,
\end{enumerate}
then every biderivation on $R$ is a sum of an extremal biderivation and an
inner biderivation.
\end{theorem}

Upper triangular matrices over a unital ring meet the two conditions for rank
at least three (Corollary~\ref{cor:higher:uppertri}) \cite[p.~1956]{Wang2016Bideriv}.

\begin{corollary}[{\cite[Corollary~2.1, p.~1956]{Wang2016Bideriv}}]\label{cor:higher:uppertri}
Let $S$ be a unital ring and let $T_n(S)$ be the ring of all $n\times n$ upper
triangular matrices over $S$ (where $n\geq 3$). Then every biderivation of
$T_n(S)$ is the sum of an extremal biderivation and an inner biderivation.
\end{corollary}

The triangular ring is the right home for the Jordan question
\cite[p.~1417]{Liu2021JordanBid}. Liu and Liu show that, under some mild
conditions, every Jordan biderivation of a triangular ring is a biderivation,
and apply it to some upper triangular matrix rings
\cite[p.~1417]{Liu2021JordanBid}. The full statement sits in their section~2
but the local copy of the paper does not extract its body; only the abstract
is quoted \cite[p.~1417]{Liu2021JordanBid}.

The same two families describe the biderivations of an incidence algebra
\cite[p.~3]{GuanZhang2024}. Guan and Zhang take a commutative ring with unity
and a locally finite poset \cite[p.~2]{GuanZhang2024}. Their inner and extremal
biderivations have the forms of Definition~\ref{def:higher:inner-extremal}
\cite[p.~3]{GuanZhang2024}. The decomposition holds once each maximal chain has
at least three elements \cite[Theorem~4.10, p.~13]{GuanZhang2024}. When every
maximal chain is infinite, the extremal part disappears and every additive
biderivation is a sum of inner ones \cite[p.~1]{GuanZhang2024}. So the incidence
algebra keeps the pattern of the triangular ring, with the shape of the
poset in place of the shape of the matrix \cite[p.~1]{GuanZhang2024}.

For a Lie algebra the space of biderivations carries a Lie bracket of its own
\cite[p.~1]{DiBartoloLaRosa2026}. Di Bartolo and La Rosa call a bilinear
map $B$ a biderivation when it is a derivation in each argument
\cite[Definition~2.1, p.~2]{DiBartoloLaRosa2026}. They also call $B$ a
right biderivation when only the first argument obeys the rule, with no
linearity asked of the second \cite[Definition~3.1, p.~3]{DiBartoloLaRosa2026}.
The right biderivations form a Lie algebra under a bracket built from the
derivations at each point \cite[Theorem~3.5, p.~4]{DiBartoloLaRosa2026}.
The left biderivations form one as well
\cite[Theorem~3.6, p.~4]{DiBartoloLaRosa2026}. For a complete
finite-dimensional Lie algebra the right space is a tensor product of functions
with the algebra itself \cite[Corollary~3.8, p.~5]{DiBartoloLaRosa2026}.
The classical biderivations are the maps that are derivations in both
arguments, and they form the first cohomology group with coefficients in the
algebra \cite[p.~1]{DiBartoloLaRosa2026}.

The same questions are asked for Leibniz algebras
\cite[p.~1]{DiBartoloDiFattaLaRosa2025Leibniz}. Two definitions of a complete
Leibniz algebra appear in the literature, and the authors treat both
\cite[Definitions~2.6 and~2.7, p.~6]{DiBartoloDiFattaLaRosa2025Leibniz}. On such
an algebra every biderivation splits into a commutator with a linear map plus a
biderivation with values in the Leibniz kernel
\cite[Theorem~3.1, p.~7]{DiBartoloDiFattaLaRosa2025Leibniz}. For a symmetric or
a skew-symmetric biderivation the first term lands in the kernel as well
\cite[Proposition~4.1, p.~10]{DiBartoloDiFattaLaRosa2025Leibniz}.

Lie superalgebras need a definition that makes the bracket itself a
superbiderivation, and the authors give one
\cite[Definition~2.6, p.~5]{DiBartoloDiFattaLaRosa2025Super}. On a complete Lie
superalgebra every superbiderivation is a bracket with a linear map of the same
degree \cite[Theorem~3.6, p.~10]{DiBartoloDiFattaLaRosa2025Super}. The authors
compare their definition with the earlier ones and record where the two
disagree \cite[pp.~1--2]{DiBartoloDiFattaLaRosa2025Super}.

A trivial extension is the triangular matrix ring of a special shape, and
Cheraghpour and Jafari describe its Jordan super-biderivations
\cite[p.~1]{CheraghpourJafari2024}. They show that such a map is a Jordan
biderivation of the base ring together with two biadditive pieces, and they record
the two identities the pieces satisfy (Theorem~\ref{thm:higher:jordan-super-bid})
\cite[Theorem~2.6, p.~9]{CheraghpourJafari2024}.\begin{theorem}[{\cite[Theorem~2.6, p.~9]{CheraghpourJafari2024}}]\label{thm:higher:jordan-super-bid}
Let $B$ be a Jordan super-biderivation of the trivial extension $T(R,M)$ of a ring
$R$ by a bimodule $M$. Then there are a Jordan biderivation $\delta$ of $R$, a
biadditive map $\beta\colon R\times M\to M$ that is a Jordan derivation in the
first coordinate with
$\beta(r,m\circ r')=\beta(r,m)\circ r'+m\circ\delta(r,r')$, a biadditive map
$\eta\colon M\times R\to M$ that is a Jordan derivation in the second coordinate
with $\eta(m\circ r',r)=\eta(m,r)\circ r'+m\circ\delta(r',r)$, and a biadditive map
$f\colon M\times M\to R$ that is a bimodule Jordan homomorphism in each coordinate,
such that
\[
B\bigl((r,m),(r',m')\bigr)=\bigl(\delta(r,r')+f(m,m'),\;\beta(r,m')+\eta(m,r')\bigr)
\]
for all $r,r'\in R$ and $m,m'\in M$.
\end{theorem}

So a Jordan super-biderivation of a trivial extension is a Jordan biderivation
\cite[p.~1]{CheraghpourJafari2024}. The same holds for an upper triangular matrix
ring under their conditions \cite[p.~1]{CheraghpourJafari2024}.

The chapter ends where it began, through Brešar, Martindale and Miers
\cite[p.~1952]{Wang2016Bideriv}. A biderivation of a noncommutative prime
ring is inner, so it has the form $\lambda[x,y]$
\cite[p.~1952]{Wang2016Bideriv}. On a triangular ring it is the sum of an
extremal and an inner piece \cite[p.~1956]{Wang2016Bideriv}. On the ring
$T_n(S)$ over a unital ring $S$ the sum holds for $n\geq 3$
\cite[p.~1956]{Wang2016Bideriv}. Liu and Liu add that a Jordan
biderivation of such a ring is a biderivation
\cite[p.~1417]{Liu2021JordanBid}.

Khan and Wani show that a multiplicative $\lambda$-$*$-Jordan triple higher derivation from a standard operator algebra on an infinite-dimensional Hilbert space into $B(H)$ is an additive $*$-higher derivation, and in particular is inner~\cite{WlMultiplicativeLambdaJordanTrip2021}.
